\documentclass[12pt,a4paper]{article}

\usepackage{ctex}
\usepackage{amsmath}
\usepackage{amssymb}
\usepackage{geometry}
\geometry{left=2.5cm,right=2.5cm,top=3cm,bottom=3cm}

\title{多元一次方程组求解——从消元法到克拉默法则}
\date{2026年2月}

\begin{document}

\maketitle

\tableofcontents

\section{问题引入：更多未知数怎么办？}

我们在初中学习过二元一次方程组，会用消元法或代入法求解。例如：

\[
\begin{cases}
x + y = 2 \\
x - y = 0
\end{cases}
\]

可以轻松求解。那三元、四元甚至更多未知数呢？

\[
\begin{cases}
x + y + z = 6 \\
x - y + z = 4 \\
x + y - z = 2
\end{cases}
\]

\textbf{核心问题}：有没有一种方法可以统一求解任意多元方程组？

\section{消元法——最普适的方法}

\subsection{消元法可以解决所有n元方程组}

\textbf{二元例子}：
\[
\begin{cases} x + y = 2 \\ x - y = 0 \end{cases}
\]

消元步骤：
\begin{enumerate}
    \item 第一个方程减去第二个方程：$(x+y) - (x-y) = 2-0$
    \item 得到 $2y = 2$，所以 $y = 1$
    \item 代入第一个方程：$x + 1 = 2$，所以 $x = 1$
\end{enumerate}

\textbf{三元例子}：
\[
\begin{cases}
x + y + z = 6 \\
x - y + z = 4 \\
x + y - z = 2
\end{cases}
\]

消元步骤：
\begin{enumerate}
    \item (1)-(2)：$2y = 2$，所以 $y = 1$
    \item (1)-(3)：$2z = 4$，所以 $z = 2$
    \item 代入(1)：$x + 1 + 2 = 6$，所以 $x = 3$
\end{enumerate}

\textbf{关键发现}：消元法可以推广到任意多元方程组！

\subsection{消元法可以判断解的情况}

在消元过程中：
\begin{itemize}
    \item 如果出现 $0 = 0$ → 可能有无穷多解
    \item 如果出现 $0 = 1$（矛盾）→ 无解
    \item 如果正好消元到只剩一个未知数 → 唯一解
\end{itemize}

\section{消元法的矩阵表示}

\subsection{从方程组到矩阵}

方程组：
\[
\begin{cases}
a_{11}x_1 + a_{12}x_2 + a_{13}x_3 = b_1 \\
a_{21}x_1 + a_{22}x_2 + a_{23}x_3 = b_2 \\
a_{31}x_1 + a_{32}x_2 + a_{33}x_3 = b_3
\end{cases}
\]

\textbf{写成矩阵形式}：
\[
A = \begin{pmatrix}
a_{11} & a_{12} & a_{13} \\
a_{21} & a_{22} & a_{23} \\
a_{31} & a_{32} & a_{33}
\end{pmatrix}, \quad
x = \begin{pmatrix} x_1 \\ x_2 \\ x_3 \end{pmatrix}, \quad
b = \begin{pmatrix} b_1 \\ b_2 \\ b_3 \end{pmatrix}
\]

简记为：$Ax = b$

其中 $A$ 称为系数矩阵，$[A|b]$ 称为增广矩阵。

\subsection{消元法对应初等行变换}

消元法的每一步，对应矩阵的行操作：

\begin{center}
\begin{tabular}{|c|c|}
\hline
消元法操作 & 矩阵操作 \\
\hline
方程乘以常数 & 行乘以常数 \\
\hline
方程相加减 & 行相加减 \\
\hline
交换两个方程 & 交换两行 \\
\hline
\end{tabular}
\end{center}

这些操作叫做\textbf{初等行变换}。

\section{初等行变换的性质}

\subsection{初等行变换不改变方程组的解}

\textbf{核心性质}：

对方程组做初等行变换，\textbf{不会改变方程组的解}！

\textbf{直观理解}：
\begin{itemize}
    \item 交换两个方程顺序 → 解不变
    \item 方程乘以非零常数 → 解不变
    \item 方程相加 → 解不变（等量替换）
\end{itemize}

\subsection{初等行变换与行列式}

对于 $2 \times 2$ 矩阵，定义\textbf{行列式}：
\[
\det\begin{pmatrix} a & b \\ c & d \end{pmatrix} = ad - bc
\]

记作：
\[
\begin{vmatrix} a & b \\ c & d \end{vmatrix}
\]

初等行变换对行列式的影响：

\begin{center}
\begin{tabular}{|c|c|}
\hline
初等行变换 & 行列式的变化 \\
\hline
交换两行 & 变号（乘以 -1） \\
\hline
行乘以常数 $k$ & 行列式乘以 $k$ \\
\hline
行乘以 $k$ 加到另一行 & \textbf{不变} \\
\hline
\end{tabular}
\end{center}

\textbf{重要}：第三种变换（倍加变换）\textbf{不改变行列式}！

这正是消元法的核心操作！

\section{对角矩阵情况——行列式的计算}

\subsection{什么是对角矩阵？}

对角矩阵：只有主对角线上有非零元素
\[
D = \begin{pmatrix}
d_1 & & \\
& d_2 & \\
& & \ddots \\
& & & d_n
\end{pmatrix}
\]

\subsection{对角矩阵的行列式}

\[
\det(D) = d_1 \cdot d_2 \cdots d_n
\]

（主对角线元素的乘积）

\subsection{对角矩阵方程组的解}

对于对角方程组：
\[
\begin{cases}
d_1 x_1 = b_1 \\
d_2 x_2 = b_2 \\
\vdots \\
d_n x_n = b_n
\end{cases}
\]

\textbf{解}：
\[
x_i = \frac{b_i}{d_i}
\]

如果某个 $d_i = 0$：
\begin{itemize}
    \item 对应 $b_i = 0$ → 该未知数任意（无穷多解）
    \item 对应 $b_i \neq 0$ → 矛盾（无解）
\end{itemize}

\textbf{规律}：解的存在取决于对角元素是否为零。

\section{行列式为零的含义——秩的概念}

\subsection{行列式为零意味着什么？}

\[
\det(A) = 0 \iff \text{矩阵的"有效信息"不足}
\]

\textbf{几何解释}：行列式为零 → 面积/体积缩放为0 → 矩阵"压扁"了空间

\subsection{从方程组角度看：秩}

\textbf{问题}：$n$ 个方程中，有多少是"有意义的"？

\textbf{例子}：
\[
\begin{cases}
x + y = 2 \\
2x + 2y = 4
\end{cases}
\]

第二个方程是第一个的 2 倍，\textbf{没有提供新信息}！

\textbf{定义}：方程组中\textbf{不能由其他方程线性组合得到}的方程个数，叫做\textbf{秩}（rank）

记作 $\text{rank}(A)$

\begin{center}
\begin{tabular}{|c|c|c|}
\hline
情况 & 秩 & 解的情况 \\
\hline
秩 = 未知数个数 & $\text{rank}(A) = n$ & 唯一解 \\
\hline
秩 < 未知数个数 & $\text{rank}(A) < n$ & 无穷多解 \\
\hline
增广矩阵秩 > 系数矩阵秩 & $\text{rank}[A|b] > \text{rank}(A)$ & 无解 \\
\hline
\end{tabular}
\end{center}

\section{克拉默法则}

\subsection{上三角矩阵与消元法}

我们已经知道：
\begin{enumerate}
    \item \textbf{消元法}：通过初等行变换将矩阵化为上三角形式
    \item \textbf{解不变}：初等行变换不改变方程组的解
    \item \textbf{上三角方程组的解}：从最后一个方程开始，依次向上回代
\end{enumerate}

设原方程组为 $Ax = b$，通过初等行变换得到上三角矩阵 $U$：

\[
\begin{pmatrix}
u_{11} & u_{12} & \cdots & u_{1n} \\
0 & u_{22} & \cdots & u_{2n} \\
\vdots & \vdots & \ddots & \vdots \\
0 & 0 & \cdots & u_{nn}
\end{pmatrix}
\begin{pmatrix}
x_1 \\ x_2 \\ \vdots \\ x_n
\end{pmatrix}
=
\begin{pmatrix}
b'_1 \\ b'_2 \\ \vdots \\ b'_n
\end{pmatrix}
\]

由于初等行变换\textbf{不改变方程组的解}，所以 $Ax=b$ 和 $Ux=b'$ 的解相同。

当 $\det(A) \neq 0$ 时，上三角矩阵 $U$ 的对角元素 $u_{11}, u_{22}, ..., u_{nn}$ 都不为零，此时方程组有唯一解。

\subsection{克拉默法则（线性列变换证明）}

\begin{center}
\fbox{
\begin{minipage}{0.9\textwidth}
\textbf{克拉默法则（Cramer's Rule）}：

对于 $n$ 元线性方程组 $Ax = b$

如果 $\det(A) \neq 0$（即 $A$ 可逆），则方程组有唯一解，且
\[
x_i = \frac{\det(A_i)}{\det(A)}
\]

其中 $A_i$ 是把 $A$ 的第 $i$ 列换成 $b$ 得到的矩阵。
\end{minipage}
}
\end{center}

\textbf{证明（利用线性列变换）}：

设 $A = (a_1, a_2, ..., a_n)$，其中 $a_j$ 是第 $j$ 列向量。

方程 $Ax = b$ 可以写成：
\[
x_1 a_1 + x_2 a_2 + \cdots + x_n a_n = b
\]

这意味着 $b$ 可以表示为 $A$ 的列向量的线性组合，系数就是 $x_1, x_2, ..., x_n$。

现在考虑矩阵 $A_i = (a_1, ..., a_{i-1}, b, a_{i+1}, ..., a_n)$，即把 $A$ 的第 $i$ 列换成 $b$。

对 $A_i$ 按第 $i$ 列展开。将 $b = x_1 a_1 + x_2 a_2 + ... + x_n a_n$ 代入：

根据行列式对列的线性性质：
\[
\det(A_i) = x_1 \det(a_1, ..., a_{i-1}, a_1, a_{i+1}, ...) + ... + x_i \det(a_1, ..., a_{i-1}, a_i, a_{i+1}, ...) + ... + x_n \det(a_1, ..., a_{i-1}, a_n, a_{i+1}, ...)
\]

注意到除了第 $i$ 项外，其他项都有两列相同（例如第一项同时包含 $a_1$ 和 $a_1$），行列式为 0。

因此：
\[
\det(A_i) = x_i \cdot \det(A) = x_i \cdot \det(A)
\]

所以：
\[
x_i = \frac{\det(A_i)}{\det(A)}
\]

这就是克拉默法则！\hfill $\square$

\textbf{说明}：这个证明利用了行列式对列的线性性质，本质上是对方程 $Ax = b$ 进行列变换，与消元法（行变换）角度不同，但结果是等价的。

\section{求证检验}

验证消元法求解的结果与克拉默法则一致。

当 $\det(A) = 0$ 时，用秩判断解的情况。

\section{结论}

\begin{enumerate}
    \item \textbf{通用求解方法}：消元法（初等行变换）
    \item \textbf{矩阵表示}：$Ax = b$
    \item \textbf{行列式性质}：初等行变换（倍加）不改变行列式
    \item \textbf{解的判断}：$\text{rank}(A)$ vs $\text{rank}[A|b]$
    \item \textbf{公式解}：克拉默法则（$\det(A) \neq 0$ 时）
\end{enumerate}

\section{后续思考}

\begin{itemize}
    \item 消元法的高效实现（LU分解）
    \item 逆矩阵与方程组求解
    \item 特征值与特征向量
\end{itemize}

\end{document}
